\subsection{热光输入：由生成函数求完整交换分布}
\label{sec:electron-thermal-exchange-analytic}

相干光的平均光子数确定以后，仍有许多不同光态具有同样的平均能量。热光就是其中一种。它没有固定光学相位，光子数也遵循与相干态不同的分布。为了求这些差别怎样传给电子，保持上一节的单模低反冲交换算符 $S_Q$ 不变，仍取电子初态 $|0\rangle_e$，只把光场输入换为热态
\begin{equation*}
 \rho_{\rm th}=\sum_{n=0}^\infty
 \frac{\bar n_{\rm th}^{\,n}}{(1+\bar n_{\rm th})^{n+1}}|n\rangle\langle n|,
 \qquad \bar n_{\rm th}=\frac1{\exp[\hbar\omega/(k_{\rm B}T_{\rm ph})]-1}.
\end{equation*}
$T_{\rm ph}$ 是该模式的温度，$\bar n_{\rm th}$ 是无量纲平均光子数。热态在光子数基中为对角混合，没有确定的光学参考相位。热光与相干光对自由电子谱的不同作用可由同一 QPINEM 模型比较\cite{Dahan2021Science}。

出射电子可以占据许多边带。求这组概率时，可以利用 Fourier 级数把所有 $P_\ell$ 合成一个函数，称为概率分布的\emph{特征函数}：
\begin{equation*}
 \chi_e(\vartheta)=\sum_{\ell\in\mathbb Z}P_\ell\ee^{\ii\ell\vartheta}
 =\operatorname{Tr}[\rho_e^{\rm out}\ee^{\ii\vartheta\hat\Lambda_e}],
 \qquad \hat\Lambda_e|\ell\rangle=\ell|\ell\rangle.
\end{equation*}
其中 $\hat\Lambda_e$ 是阶数算符，$\vartheta$ 是无量纲 Fourier 变量。它只是用来整理概率，不代表额外的电子动力学。因为 $\partial_\vartheta\ee^{\ii\ell\vartheta}=\ii\ell\ee^{\ii\ell\vartheta}$，在原点求导便有
\[
\langle\ell\rangle=-\ii\chi_e'(0),\qquad
\langle\ell^2\rangle=-\chi_e''(0),\qquad \chi_e(0)=1.
\]
对数的导数则自动去掉低阶矩的乘积：
\[
-\ii(\ln\chi_e)'(0)=\langle\ell\rangle,\qquad
-(\ln\chi_e)''(0)=\langle\ell^2\rangle-\langle\ell\rangle^2.
\]
这些对数导数称为\emph{累积量}，前两阶就是均值和方差。用本节末尾的算符代数求热态平均，得到
\begin{equation}
 \chi_e(\vartheta)=
 \exp\!\left\{\lambda_+(\ee^{\ii\vartheta}-1)
             +\lambda_-(\ee^{-\ii\vartheta}-1)\right\},
 \qquad
 \lambda_+=r_Q\bar n_{\rm th},\quad
 \lambda_-=r_Q(1+\bar n_{\rm th}).
 \label{eq:electron-thermal-characteristic}
\end{equation}
该结果对本交换模型中任意有限 $r_Q=|\gq|^2$ 成立，不是只保留单量子的微扰式。$\lambda_\pm$ 是无量纲统计参数；形式上等同于两个独立 Poisson 变量 $N_\pm$ 的均值，净电子阶数为 $\ell=N_+-N_-$。这是一种\emph{末态分布的等价表示}，不自动意味着真实连续时间动力学是两个无记忆跳跃过程。

直接对 $\ln\chi_e$ 求导得
\begin{equation}
 \langle\ell\rangle=\lambda_+-\lambda_-=-r_Q,\qquad
 \operatorname{Var}\ell=\lambda_++\lambda_-=r_Q(2\bar n_{\rm th}+1).
 \label{eq:electron-thermal-moments}
\end{equation}
热光增强增益和损失两侧的涨落，但净损失偏置仍来自真空项。取 $\bar n_{\rm th}=|\alpha_{\rm coh}|^2$，热光与前节相干光的均值和方差完全相同；要区分两者，必须读更高统计矩、谱形或相干，不能只读展宽。

\begin{exbox}[解析恢复热交换概率并比较四阶累积量]
从特征函数求 $P_\ell$。对 $\ell\ge0$，条件 $N_+-N_-=\ell$ 等价于 $N_+=m+\ell,N_-=m$，故
\begin{equation*}
 P_\ell=\ee^{-(\lambda_++\lambda_-)}
 \sum_{m=0}^\infty\frac{\lambda_+^{m+\ell}\lambda_-^m}{(m+\ell)!m!}.
\end{equation*}
用第一类修正 Bessel 函数的定义
$I_j(2x)=\sum_{m\ge0}x^{2m+j}/[m!(m+j)!]$，并对负 $\ell$ 交换两参数，得到
\begin{equation*}
 P_\ell=\ee^{-(\lambda_++\lambda_-)}
 \left(\frac{\lambda_+}{\lambda_-}\right)^{\ell/2}
 I_{|\ell|}(2\sqrt{\lambda_+\lambda_-}),\qquad \bar n_{\rm th}>0.
\end{equation*}
它称为两个 Poisson 变量之差的 Skellam 分布。零温不能直接把零代入负幂；应回到求和式或特征函数，得到
$P_{-m}=\ee^{-r_Q}r_Q^m/m!$（$m\ge0$），其余增益概率为零。对任意正温度还可直接相除：
\begin{equation*}
 \frac{P_{+\ell}}{P_{-\ell}}
 =\left(\frac{\bar n_{\rm th}}{1+\bar n_{\rm th}}\right)^\ell
 =\ee^{-\ell\hbar\omega/(k_{\rm B}T_{\rm ph})},\qquad\ell>0.
\end{equation*}
这给出了不依赖总体耦合强度的增益--损失比。

再定义四阶累积量 $\mathcal C_4=\langle(\ell-\langle\ell\rangle)^4\rangle-3(\operatorname{Var}\ell)^2$。它是对高斯统计之外结构的一个度量。由\eqnrefs{eq:electron-thermal-characteristic}，
\begin{equation*}
 \mathcal C_4^{\rm(th)}=r_Q(2\bar n_{\rm th}+1).
\end{equation*}
相干光的经典 Bessel 部分具有特征函数 $J_0[4|\beta|\sin(\vartheta/2)]$，真空 Poisson 损失与它独立。将 $\ln J_0$ 展开至 $\vartheta^4$，得到
\begin{equation*}
 \mathcal C_4^{\rm(coh)}=r_Q+2|\beta|^2-6|\beta|^4.
\end{equation*}
在相同平均光子数、$|\beta|^2=r_Q\bar n_{\rm th}$ 下，两种方差相同，四阶累积量却相差 $6r_Q^2\bar n_{\rm th}^2$。这说明为什么只拟合一个高斯宽度会丢失光子统计信息。
\end{exbox}

\begin{derivation}[用交换算符代数完成热态偏迹]
定义 $B=\hat b^\dagger\hat a$。双边平移给 $[B,B^\dagger]=I$，因此可定义位移形式
$D_B(\zeta)=\exp(\zeta B^\dagger-\zeta^*B)$，并写 $S_Q=D_B(-\gq^*)$。
电子计数旋转满足
\begin{equation*}
 \ee^{\ii\vartheta\hat\Lambda_e}S_Q\ee^{-\ii\vartheta\hat\Lambda_e}
 =D_B(-\gq^*\ee^{-\ii\vartheta}).
\end{equation*}
由于初态 $|0\rangle_e$ 的阶数为零，$\chi_e$ 可改写为初态上 $S_Q^\dagger S_Q(\vartheta)$ 的期望。使用
\begin{equation*}
 D_B(\alpha)D_B(\beta)
 =\ee^{(\alpha\beta^*-\alpha^*\beta)/2}D_B(\alpha+\beta)
\end{equation*}
得到
\begin{equation*}
 \chi_e(\vartheta)=\ee^{-\ii r_Q\sin\vartheta}
 \left\langle D_B[-\gq^*(\ee^{-\ii\vartheta}-1)]\right\rangle_{0,\rm th}.
\end{equation*}
接下来显式求平均。正规排序给
$D_B(\zeta)=\ee^{-|\zeta|^2/2}\ee^{\zeta B^\dagger}\ee^{-\zeta^*B}$。
在 $|0,n\rangle$ 上，只有同次数的 $B^\dagger$ 与 $B$ 能返回原电子阶数，因此
\begin{equation*}
 \langle0,n|D_B(\zeta)|0,n\rangle
 =\ee^{-|\zeta|^2/2}\sum_{k=0}^n
 \frac{(-|\zeta|^2)^k}{(k!)^2}\frac{n!}{(n-k)!}
 =\ee^{-|\zeta|^2/2}L_n(|\zeta|^2).
\end{equation*}
$L_n$ 是普通 Laguerre 多项式，上一式的有限和也可作为其定义。令 $q_{\rm th}=\bar n_{\rm th}/(1+\bar n_{\rm th})$，用几何级数加权后的恒等式
\begin{equation*}
 (1-q_{\rm th})\sum_{n=0}^\infty q_{\rm th}^nL_n(x)
 =\exp[-\bar n_{\rm th}x]
\end{equation*}
这个恒等式可由几何级数逐项求导验证：对 $\sum_{n\ge0}q^n=(1-q)^{-1}$ 求 $k$ 阶导数并乘 $q^k$，有
\begin{equation*}
 (1-q)\sum_{n=k}^{\infty}q^n\frac{n!}{(n-k)!}
 =k!\left(\frac{q}{1-q}\right)^k.
\end{equation*}
将它代入 $L_n$ 的有限和，剩下 $\sum_k[-xq/(1-q)]^k/k!$，正是指数级数。于是得到 $\langle D_B(\zeta)\rangle_{0,\rm th}
=\exp[-(\bar n_{\rm th}+1/2)|\zeta|^2]$。
最后代入 $|\zeta|^2=2r_Q(1-\cos\vartheta)$，
\begin{equation*}
 \ln\chi_e=-\ii r_Q\sin\vartheta
 -r_Q(2\bar n_{\rm th}+1)(1-\cos\vartheta),
\end{equation*}
按 $\ee^{\pm\ii\vartheta}$ 重组即得\eqnrefs{eq:electron-thermal-characteristic}。热态对光学相位旋转不变，初始电子又在单一阶数上；交换数守恒遂使电子约化态在阶数基中保持对角。这里的 $P_\ell$ 因而给出了完整约化电子态，而非只给其一部分矩阵元。
\end{derivation}
